\section{分层证据核验方法}
\subsection{对象、声明和允许的结论}
把一个被评价对象写为
\begin{equation}
x=(s_v,q,m,a,e,b),
\end{equation}
其中 $s_v$ 是固定版本的技能包，$q$ 为任务，$m$ 为模型，$a$ 为执行器，
$e$ 为依赖和权限环境，$b$ 为上下文、时间及费用预算。
质量分不是只由技能名称决定的常量。
对每条能力声明 $c$，保存其输入、输出、前提、排除项、来源及版本，
再按表\ref{tab:states}记录当前证据状态。
来源哈希和行号有效说明该窗口可以被回放；它们不自动证明窗口支持声明，
更不证明相应行为在当前环境中成功。

\begin{table}[t]
\caption{证据状态及其允许支持的判断。状态不是从低到高可相加的质量等级。}
\label{tab:states}
\centering\small
\begin{tabularx}{\columnwidth}{p{.24\columnwidth}X}
\toprule
状态 & 允许的解释 \\
\midrule
仅声明 & 包或模型声称具备某能力，尚未得到验证。\\
静态支持 & 在登记资源和支持的分析范围内与声明一致。\\
发现冲突 & 有版本绑定的反证；需确认整体影响。\\
产物已验收 & 在指定任务、模型、执行器及预算下通过独立验收。\\
未知 & 缺材料、超出支持范围或尚未运行，不能补成零分或安全。\\
\bottomrule
\end{tabularx}
\end{table}

现有八维锚点包括契约一致性、完整性、静态可执行性、验证质量、维护性、
来源许可、成本意识和安全说明质量。
这些维度描述可观察材料；风险门禁、任务适配和真实效用另行测量。
纯咨询技能没有脚本，不因此被判为低质量。
反之，即使一个包的文档维度很高，也不能用其平均值抵消已发现的实质接口冲突。

\subsection{完整包与静态契约}
先核对资源清单、来源提交、字节数和 SHA-256，再分析受支持文本。
脚本只经语法树解析，绝不导入或执行社区代码。
本轮契约检查限于单一字面 argparse 解析器，支持选项存在性、必填、枚举值、
简单数值类型和布尔开关。
动态解析器、子命令、复杂位置参数和无法可靠解析的示例返回未知。
这种受限性使判断可反驳，也要求同时报告未知比例，而不能只报告命中案例。

对于受支持命令，构造已知错误和无害变换，比较只查参数名的基线、扩展契约检查及
按白名单声明重建的标准库 argparse。
后两者的约束逻辑不同，但共享语法树提取与词法预处理，不能被视为完全独立的端到端真值系统。
缺陷变体用于验证该机械契约，不充当整个技能的质量标签。

\subsection{来源绑定与语义支持分离}
模型返回的引用包含登记路径、源哈希、起止行、原文和支持或反证关系。
严格定位器只接受已提供的带编号材料，检查所有字段是否吻合。
随后，系统来源绑定器可以在冻结清单内额外取回资源，检查字节与哈希，并由系统抽取有界窗口。
模型不能利用引用路径读取清单之外的本机文件。

若模型只摘录了标注行的一部分，系统保留这种差异，不把重新抽取的完整窗口倒算成原模型正确。
若引用指向仅有目录登记的文件，后续补证单列为系统取证，不能说模型先前已经看过其正文。
绑定状态与语义支持状态使用独立字段，后者在未经审查时保持未评估。
这一设计借鉴事实与引用评价\citep{P39,P40,P41}，将版本绑定、定位、语义与结论覆盖拆开，
避免把一个可计算的引用有效率误用为综合质量准确率。

\begin{figure*}[t]
\centering
\begin{tikzpicture}[box/.style={draw=blue!45!black,rounded corners=2pt,fill=blue!3,align=center,minimum height=10mm,text width=34mm},>=Stealth,font=\small]
\node[box] (v) {固定版本与资源清单\\来源、字节、哈希};
\node[box,right=8mm of v] (c) {原子声明与静态反证\\接口、资源、前提};
\node[box,right=8mm of c] (b) {模型判断与来源绑定\\保留差异和未知};
\node[box,right=8mm of b] (o) {独立产物验收\\成功、失败及成本};
\draw[->] (v)--(c);\draw[->] (c)--(b);\draw[->] (b)--(o);
\node[draw,dashed,align=center,below=7mm of b,text width=70mm,minimum height=9mm] (u) {新实现：证据判别与条件化路径；仍待外部校准};
\draw[->,dashed] (o.south)|-(u.east);\draw[->,dashed] (u.west)-|(c.south);
\end{tikzpicture}
\caption{分层证据协议。实线为分层测量，虚线连接后文新增判别器与路径方法；外部效度仍待检验。来源绑定不替代语义判断；产物验收不等于跨任务的永久能力认证。}
\label{fig:protocol}
\end{figure*}

\subsection{独立产物验收与条件化效用}
对于任务 $q$ 的输出产物 $o$，令 $C_q$ 是预先定义的必需检查集合。
完整任务成功定义为
\begin{equation}
Y_q(o)=\prod_{k\in C_q}\mathbf{1}[k(o)=\mathrm{pass}],
\label{eq:success}
\end{equation}
而非各子检查通过比例的平均。
数值类型也属于契约，例如布尔真不能替代整数 1，浮点 8.0 不能替代要求的整数 8。
范围外或前提缺失的正确处置单独计量，不能与范围内的产物成功混成一个比例。

技能的任务条件化贡献为
\begin{equation}
\Delta_s=\frac{1}{H}\sum_{h=1}^{H}\frac{1}{|T_h|}
\sum_{t\in T_h}\bigl[Y_t(s)-Y_t(\varnothing)\bigr],
\label{eq:utility}
\end{equation}
其中 $h$ 为任务条件簇，$T_h$ 为该簇的输入与重复组合。
比较必须固定基础模型、必需任务说明、输出上限和可用工具，同时记录技能占用的上下文。
本文没有把无 Skill 组填充成无意义等长文本。
额外的无关技能组和仅操作契约组用于诊断上下文干扰及新增说明的贡献；
观测 token 不相等时如实披露，不宣称精确预算匹配。

经验能力边界是被测 $(q,m,a,e,b)$ 条件上的结果集合。
未观测到成功不证明数学意义的不可能，少数成功也不足以认证整个输入区间。
本研究不设未经校准的通用“90\% 即可自主使用”阈值，也不提出共形安全保证。

\subsection{技能网络的集合测量}
对任务 $q$ 的参考必要技能集 $G_q$ 与选择集 $S_q$，分别定义
\begin{align}
R@k(q)&=\frac{|G_q\cap S_q|}{|G_q|},\\
F(q)&=\mathbf{1}[G_q\subseteq S_q].
\end{align}
前者衡量平均参考覆盖，后者要求整组必要能力齐备。
二者都受参考标签遗漏等价技能的影响，并不是实际任务成功率。
将选择空间限制到固定前 20 个候选、固定取 5 个后，比较原顺序、MMR、DPP、
查询残差 DPP 和查询覆盖，检验“增加多样性总会更好”这一假设。
本文使用词法特征与名次相关性代理；没有继承 DSR 的训练编码器和原论文效果\citep{P05,P38}。
